\documentclass[11pt,letterpaper]{article}
\usepackage{fontspec}
\usepackage{tocloft}
\usepackage{multicol}
\usepackage{nicefrac}
\usepackage{url}
\usepackage[left=1.3in,right=1.3in,top=1in,bottom=1in]{geometry}
\setmainfont[Scale=1.1,AutoFakeBold=1.5,AutoFakeSlant=0.3]{Doves Type}

% Prevent page breaks within list items
\widowpenalty=10000
\clubpenalty=10000
\interlinepenalty=500

\title{Pommes Boulang\`ere •}
\author{}
\date{}

\begin{document}

\maketitle
\thispagestyle{empty}

\begin{quote}\textit{A bed of long-caramelized onions carries overlapping layers of thinly sliced potato, each one seasoned and moistened with clarified butter and beef stock. Baked low under a parchment cartouche and a weighted plate, the potatoes steam and absorb the stock without browning; the cover comes off and the oven rises only to \textit{265°F}, where repeated spoonfuls of stock and butter glaze the rim to a chewy, lacquered amber while the center stays tender. Thyme goes on in the last five minutes. Serves 4--6 as a side.}\end{quote}

\section*{Ingredients}
\setlength{\columnsep}{20pt}
\begin{multicols}{2}
\noindent
    Yellow onions \dotfill 3~large (about 1\nicefrac{1}{2}~lb.) \\
    Neutral oil \dotfill 2~Tbsp. \\
    Kosher salt \dotfill 1~tsp. \\
    Starchy potatoes (Russet or Maris Piper) \dotfill 2~lb. \\
    Unsalted butter \dotfill 8~oz. (2~sticks) \\

    \columnbreak
    \noindent
    Beef base (Better Than Bouillon) \dotfill 1~Tbsp. \\
    Water \dotfill 3~cups \\
    Black pepper, freshly ground \dotfill \nicefrac{3}{4}~tsp. \\
    Fresh thyme, leaves picked \dotfill 2~Tbsp. \\
\end{multicols}

\section*{Directions}

\noindent
Preheat oven to \textit{250°F} ---
Halve and slice \textbf{onions} about \nicefrac{1}{8}~inch thick ---
Pick \textbf{thyme} leaves from the stems into \textit{Small Bowl~\#1} ---
Combine the remaining \nicefrac{1}{2}~tsp. \textbf{salt} with the \nicefrac{3}{4}~tsp. \textbf{black pepper} in \textit{Small Bowl~\#2} (layer seasoning) ---
Stir 1~Tbsp. \textbf{beef base} into 3~cups hot \textbf{water} in a 4-cup measuring cup until dissolved; keep hot ---
Slice \textbf{potatoes} \nicefrac{1}{8}~inch thick on a mandoline with the guard on; keep the larger slices apart from the smaller ones ---
Melt \textbf{butter} in a microwave-safe jug in \textit{20-second bursts until fully liquid}, about \textit{1--2~minutes}, and leave it to stand so the clear fat rises; do not stir it back together

\begin{enumerate}
    \item Caramelize the \textbf{onions} in a 9-inch oven-safe skillet over medium-low heat: add \textbf{neutral oil} and the sliced \textbf{onions} with \nicefrac{1}{2}~tsp. of the \textbf{kosher salt}, and stir to coat. Cook for \textit{30--40~minutes}, stirring every few minutes and scraping the base of the pan, until they collapse to a soft, jammy mass the colour of dark honey. Halfway through they will look wet and translucent---water leaving, not sugar browning---so keep going. Keep the heat low so that no edge blackens, and if they start to catch, drop the heat and add a splash of water. They are ready when a spoon dragged across the pan leaves a trail that fills slowly and the smell is sweet rather than sharp. Turn off the heat.

    \item Taste the onion bed and adjust. Drizzle \textit{2~Tbsp.} of the clear \textbf{clarified butter} over the \textbf{onions}, leaving the milk solids behind in the jug, and add \textit{\nicefrac{1}{2}~cup} of the hot \textbf{beef stock}. Stir to combine and spread level across the base of the pan. Taste a spoonful: it should be savoury and gently salted, not salty---this bed is the deepest flavour in the dish and cannot be corrected later.

    \item Lay the first layer of \textbf{potatoes}. Arrange the larger slices around the rim and work inward, overlapping each slice by about a third and finishing in the center with the smallest slices and least overlap; the thicker thatch at the rim is what survives the bake and caramelizes as one piece. Ladle \textit{\nicefrac{1}{3}~cup} of the hot \textbf{stock} over the layer, season with a pinch of \textbf{layer seasoning} (\textit{Small Bowl~\#2}), and drizzle \textit{1~Tbsp.} of \textbf{clarified butter}.

    \item Repeat the layering until the potatoes are used up, five to seven layers in all.
\begin{enumerate}
    \item Start each new layer from the opposite side of the pan, so the slices bridge the seams underneath and no gap runs straight down through the stack.
    \item Scatter a third of the \textbf{thyme} (\textit{Small Bowl~\#1}) over the middle layer only; the leaves on the surface would scorch long before the bake ends.
    \item Season and moisten every layer the same way, but hold back the last \textit{3~Tbsp.} of \textbf{butter} and the last \nicefrac{1}{2}~cup of \textbf{stock} for the glaze.
    \item Finish with a plain layer, cut faces up, giving the rim the heaviest share of the butter and stock: it cooks fastest and it is the part everyone reaches for.
\end{enumerate}

    \item Cover the potatoes with a cartouche. Fold a sheet of parchment in quarters, lay it on the pan, cut around the rim, then snip a \nicefrac{1}{2}~inch hole in the center. Unfold it over the potatoes and press the overhang down against the inside wall all the way around so the paper seals to the pan---the hole is the only place steam can leave, and the seal is what keeps the layers from drying.
\begin{enumerate}
    \item Set a heatproof plate on top of the parchment as a weight; it holds the layers flat while the potatoes soften and the stock moves through them.
    \item Bake at \textit{250°F} for \textit{45~minutes}, sliding the plate out and turning it over once at the halfway point so the stack flattens evenly.
\end{enumerate}

    \item Remove the plate and set it aside for good---the layers have set and hold their line without it. Raise the oven to \textit{265°F} and bake for a further \textit{25~minutes} with the cartouche still on. When it comes off, the stock should be bubbling at the rim and the paper should be stained and slightly shrunken but not browned.

    \item Remove the cartouche and glaze the top for the rest of the bake. All the fat and stock have risen and the pan is still wet, so add nothing yet: tilt the pan and spoon the pooled juices over the top every \textit{8--10~minutes} for \textit{20--25~minutes}, until the rim is a deep lacquered amber and the liquid has reduced to a film coating the potatoes. Dot any pale patch with reserved \textbf{butter}, and add reserved \textbf{stock} only where the pan looks dry.
\begin{enumerate}
    \item Do not be tempted to raise the oven to brown them faster. Above \textit{265°F} the top layer shrivels and dries before the sugars darken; the glaze builds from repeated thin coats, not from heat.
    \item The potatoes are done when a knife slides through the center with no resistance at all and the rim is caramelized but never black.
\end{enumerate}

    \item Finish with \textbf{thyme}. Spoon the last of the juices into the dips on the surface, scatter the remaining \textbf{thyme} (\textit{Small Bowl~\#1}), and return to the oven for \textit{5~minutes}---any longer and the leaves burn. Rest the pan \textit{5~minutes} off the heat, then scoop down through the layers at the table so every portion carries crust, tender potato, and onion.
\end{enumerate}

\newpage

{\footnotesize
\setlength{\columnsep}{20pt}
\setlength{\multicolsep}{6pt}
\begin{multicols}{2}
\setlength{\parindent}{0pt}
\setlength{\parskip}{2pt}
\setlength{\itemsep}{0pt}
\setlength{\parsep}{0pt}

\subsection*{Yield}
\begin{itemize}
    \item Serves 4--6 as a side, from one 9-inch pan
\end{itemize}

\subsection*{Equipment Required}
\begin{itemize}
    \item 9-inch oven-safe non-stick or cast iron skillet, about \textit{2\nicefrac{1}{2}--3~inches} deep
    \item Mandoline with hand guard (or a sharp chef's knife and patience)
    \item Parchment paper for the cartouche
    \item Heatproof dinner plate or metal cake pan, 8--9~inches across, as a weight
    \item Microwave-safe jug or bowl (1~quart) for the \textbf{butter}
    \item 4-cup measuring cup for the \textbf{stock}
    \item Ladle and spoon for basting
    \item Oven mitts, tongs, cutting board, measuring spoons
    \item Small prep bowls (2)
\end{itemize}

\subsection*{Mise en Place}
\begin{itemize}
    \item Small Bowl \#1 --- picked \textbf{thyme} leaves
    \item Small Bowl \#2 --- \nicefrac{1}{2}~tsp. \textbf{salt} and \nicefrac{3}{4}~tsp. \textbf{black pepper} combined (layer seasoning)
    \item 4-cup measuring cup --- 1~Tbsp. \textbf{beef base} dissolved in 3~cups hot \textbf{water}; keep hot
    \item Microwave-safe jug --- \textbf{butter}, melted in \textit{20-second bursts}, clear fat left on top of the milk solids
    \item \textbf{Onions} sliced first: they take the longest and cook while everything else is prepped
    \item \textbf{Potatoes} sliced last, just before layering, with larger and smaller slices kept apart
\end{itemize}

\subsection*{Ingredient Tips}
\begin{itemize}
    \item \textbf{Potatoes}: a floury, starchy variety---Russet, Maris Piper, King Edward, Agria. Waxy potatoes stay firm and refuse to drink the stock
    \item \textbf{Beef base}: Better Than Bouillon is heavily salted, which is why the \textbf{salt} here is a single teaspoon and the first layer gets tasted before the rest is built
    \item \textbf{Butter}: unsalted, so the \textbf{beef base} controls the seasoning. The \textit{8~oz.} melts down to roughly \textit{12~Tbsp.} of clear fat, which is what the layers are built on. Leaving the milk solids in the jug is fine---they would only scorch once the cover comes off
    \item \textbf{Stock}: real beef stock is better than base; if you use it, expect to add a little more \textbf{salt}
    \item \textbf{Thyme}: leaves picked from the stems. Stems left in will burn to black threads at \textit{265°F}
\end{itemize}

\subsection*{Preparation Tips}
\begin{itemize}
    \item Uniform \nicefrac{1}{8}-inch slices are the whole technique. Thick slices stay raw next to thin ones that have already dissolved
    \item Slice on a mandoline with the guard on, and use the palm of your hand against the guard, never fingertips on the blade
    \item Overlap lightly. A heavy overlap doubles the layers at the seam and leaves them undercooked in the middle of the pan
    \item The cartouche must seal against the pan wall. If steam escapes at the edge, the top layers dry out and curl instead of steaming
    \item Every layer wants salt, pepper, butter, and stock. Seasoning only the top is the most common failure of this dish
    \item Judge the finish by colour and by feel, not by the clock
\end{itemize}

\subsection*{Make Ahead \& Storage}
\begin{itemize}
    \item Caramelize the \textbf{onions} up to \textit{3~days} ahead and refrigerate; the bed can be built directly from cold, allowing \textit{10~minutes} more in the first bake
    \item Assemble the full dish, cover, and refrigerate up to \textit{24~hours}. Bake from cold and add \textit{15~minutes} to the covered stage
    \item Leftovers keep \textit{3~days} refrigerated. Reheat at \textit{350°F} for \textit{20--25~minutes} with a splash of \textbf{stock} over the top to bring the glaze back
    \item Do not freeze: the layers break down and the texture is lost
\end{itemize}

\subsection*{Serving Suggestions}
\begin{itemize}
    \item Roast lamb, beef, or chicken. The pan juices and the onion bed are the sauce; gravy is unnecessary
    \item A bitter green---watercress, chicory, or buttered greens---cuts the butter
    \item Serve in the pan. The caramelized rim is the point of the dish and it softens within \textit{15~minutes} of coming out
\end{itemize}

\subsection*{Source}
\begin{itemize}
    \item Adapted from Will Murray (Fallow, London), via Fallow: \url{https://www.youtube.com/watch?v=-Z3P_NujFos}
\end{itemize}

\end{multicols}
}

\end{document}
